\chapter{1987年全国卷（文）}

\section{单选题}

\begin{problem}
设 \(S\)，\(T\) 是两个非空集合，且 \(S \not\subseteq T\)，\(T \not\subseteq S\)，令 \(X = S \cap T\)，那么 \(S \cup X\) 等于（\quad）

\choices
    {\(X\)}
    {\(T\)}
    {\(\varnothing\)}
    {\(S\)}
\end{problem}

\begin{answer}
D
\end{answer}

\begin{solution}
由 \(X=S\cap T\) 可知，\(X\subseteq S\). 集合与其子集的并仍等于这个集合，因此 \(S\cup X=S\)，故选 D.
\end{solution}

\begin{problem}
设椭圆方程为 \(\frac{x^2}{a^2} + \frac{y^2}{b^2} = 1\)（\(a > b > 0\)），令 \(c = \sqrt{a^2 - b^2}\)，那么它的准线方程为（\quad）

\choices
    {\(y = \pm \frac{a^2}{c}\)}
    {\(y = \pm \frac{b^2}{c}\)}
    {\(x = \pm \frac{a^2}{c}\)}
    {\(x = \pm \frac{b^2}{c}\)}
\end{problem}

\begin{answer}
C
\end{answer}

\begin{solution}
因为 \(a>b>0\)，分母较大的 \(a^2\) 对应长轴，所以椭圆的长轴在 \(x\) 轴上，焦点为 \(F_1(c,0)\) 和 \(F_2(-c,0)\). 其离心率为 \(e=\frac ca\). 椭圆的准线垂直于长轴，且标准椭圆的两条准线分别为 \(x=\pm\frac ae\). 代入 \(e=\frac ca\)，得 \(x=\pm\frac{a^2}{c}\)，故选 C.
\end{solution}

\begin{problem}
设 \(\log_3 4 \cdot \log_4 8 \cdot \log_8 m = \log_4 16\)，那么 \(m\) 等于（\quad）

\choices
    {\(\frac{9}{2}\)}
    {\(9\)}
    {\(18\)}
    {\(27\)}
\end{problem}

\begin{answer}
B
\end{answer}

\begin{solution}
由对数换底公式，\(\log_3 4\cdot\log_4 8=\log_3 8\)，从而原等式左边为 \(\log_3 8\cdot\log_8m=\log_3m\). 右边 \(\log_4 16=2\)，所以 \(\log_3m=2\)，解得 \(m=3^2=9\). 该值满足对数真数 \(m>0\)，故选 B.
\end{solution}

\begin{problem}
复数 \(\sin 40^\circ - \i\cos 40^\circ\) 的辐角为（\quad）

\choices
    {\(40^\circ\)}
    {\(140^\circ\)}
    {\(220^\circ\)}
    {\(310^\circ\)}
\end{problem}

\begin{answer}
D
\end{answer}

\begin{solution}
利用余函数关系，\(\sin40^\circ=\cos50^\circ\)，\(\cos40^\circ=\sin50^\circ\)，因此 \(\sin40^\circ-\i\cos40^\circ=\cos50^\circ-\i\sin50^\circ\). 又 \(\cos(-50^\circ)=\cos50^\circ\)，\(\sin(-50^\circ)=-\sin50^\circ\)，所以该复数可写为 \(\cos(-50^\circ)+\i\sin(-50^\circ)\). 它的辐角为 \(-50^\circ+360^\circ k\)（\(k\in\Z\)），在给出的 \(0^\circ\) 到 \(360^\circ\) 的选项中对应 \(310^\circ\)，故选 D.
\end{solution}

\begin{problem}
二次函数 \(y = f(x)\) 的图象如图所示，那么此函数为（\quad）

\FigureLayoutDeclare{img/1987/national_liberal/q5_fig1.png}{6.0cm}{25bp 13bp 17bp 25bp}
\bitmapfigure[max width=0.44\linewidth]{img/1987/national_liberal/q5_fig1.png}

\choices
    {\(y = x^2 - 4\)}
    {\(y = 4 - x^2\)}
    {\(y = \frac{3}{4}\left(4 - x^2\right)\)}
    {\(y = \frac{3}{4}\left(2 - x^2\right)\)}
\end{problem}

\begin{answer}
C
\end{answer}

\begin{solution}
由图象可知抛物线的顶点为 \((0,3)\)，且与 \(x\) 轴交于 \((-2,0)\) 和 \((2,0)\). 因为对称轴为 \(y\) 轴，可设函数为 \(f(x)=Ax^2+3\). 代入图象上的点 \((2,0)\)，得 \(4A+3=0\)，所以 \(A=-\frac34\). 于是 \(f(x)=3-\frac34x^2=\frac34\left(4-x^2\right)\)，故选 C.
\end{solution}

\begin{problem}
在区间 \((-\infty,0)\) 上为增函数的是（\quad）

\choices
    {\(y = -\log_{\frac{1}{2}}(-x)\)}
    {\(y = \frac{x}{1 - x}\)}
    {\(y = -(x + 1)^2\)}
    {\(y = 1 + x^2\)}
\end{problem}

\begin{answer}
B
\end{answer}

\begin{solution}
在区间 \((-\infty,0)\) 内分别判断各函数的导数. 选项 A 中 \(y'=-\frac{1}{x\ln\frac12}<0\)，因为 \(x<0\) 且 \(\ln\frac12<0\)，所以它是减函数. 选项 B 中 \(y'=\frac{(1-x)+x}{(1-x)^2}=\frac{1}{(1-x)^2}>0\)，所以它在整个区间上为增函数. 选项 C 中 \(y'=-2(x+1)\)，当 \(x<-1\) 时导数为正，当 \(-1<x<0\) 时导数为负，不能在整个区间上保持增性. 选项 D 中 \(y'=2x<0\)，为减函数. 因此只有选项 B 符合，故选 B.
\end{solution}

\begin{problem}
已知平面上一点 \(P\) 在原坐标系中的坐标为 \((0,m)\)（\(m \neq 0\)），而在平移后的新坐标系中的坐标为 \((m,0)\)，那么新坐标系的原点 \(O'\) 在原坐标系中的坐标为（\quad）

\choices
    {\((-m,m)\)}
    {\((m,-m)\)}
    {\((m,m)\)}
    {\((-m,-m)\)}
\end{problem}

\begin{answer}
A
\end{answer}

\begin{solution}
设新坐标系原点 \(O'\) 在原坐标系中的坐标为 \((h,k)\). 平移只改变原点，不改变坐标轴的方向，所以任一点的“新坐标”都等于它的原坐标减去 \((h,k)\). 点 \(P\) 的两组坐标满足 \((0-h,m-k)=(m,0)\). 比较两个坐标，得到 \(h=-m\)，\(k=m\)，所以 \(O'\) 的原坐标为 \((-m,m)\)，故选 A.
\end{solution}

\begin{problem}
要得到函数 \(y = \sin \left(2x - \frac{\pi}{3}\right)\) 的图象，只需将函数 \(y = \sin 2x\) 的图象（如图）（\quad）

\bitmapfigure[max width=0.5\linewidth]{img/1987/national_liberal/q8_fig1.png}

\choices
    {向左平行移动 \(\frac{\pi}{3}\)}
    {向右平行移动 \(\frac{\pi}{3}\)}
    {向左平行移动 \(\frac{\pi}{6}\)}
    {向右平行移动 \(\frac{\pi}{6}\)}
\end{problem}

\begin{answer}
D
\end{answer}

\begin{solution}
将自变量中的式子配成 \(2\) 与平移量的形式：\(2x-\frac\pi3=2\left(x-\frac\pi6\right)\). 因此 \(y=\sin\left(2\left(x-\frac\pi6\right)\right)\) 是把 \(y=\sin2x\) 中的 \(x\) 换成 \(x-\frac\pi6\) 得到的图象，也就是向右平移 \(\frac\pi6\)，故选 D.
\end{solution}

\section{解答题}

\begin{problem}
求函数 \(y = \sin^2 2x\) 的周期.
\end{problem}

\begin{solution}
由降幂公式，\(y=\sin^2 2x=\frac{1-\cos4x}{2}\). 函数 \(\cos4x\) 的最小正周期为 \(\frac{2\pi}{4}=\frac\pi2\)，加上常数并乘以非零常数不会改变周期，因此 \(y=\sin^22x\) 的最小正周期为 \(\frac\pi2\).
\end{solution}

\begin{problem}
已知方程 \(\frac{x^2}{2 + \lambda} - \frac{y^2}{1 + \lambda} = 1\) 表示双曲线，求 \(\lambda\) 的范围.
\end{problem}

\begin{solution}
当 \(\lambda=-2\) 或 \(\lambda=-1\) 时，方程中的一个分母为零，方程没有意义，所以这两个值必须排除. 将方程写成 \(Ax^2+By^2=1\)，其中 \(A=\frac{1}{\lambda+2}\)，\(B=-\frac{1}{\lambda+1}\). 要表示非退化双曲线，\(x^2\) 与 \(y^2\) 的系数必须异号，即 \(AB<0\). 于是 \(AB=-\frac{1}{(\lambda+2)(\lambda+1)}<0\) 等价于 \((\lambda+2)(\lambda+1)>0\). 在临界点 \(-2\)，\(-1\) 两侧作符号判断，得 \(\lambda<-2\) 或 \(\lambda>-1\).
\end{solution}

\begin{problem}
若 \((1 + x)^n\) 的展开式中，\(x^3\) 的系数等于 \(x\) 的系数的 \(7\) 倍，求 \(n\).
\end{problem}

\begin{solution}
这里 \(n\) 为正整数. 若 \(n<3\)，则 \(x^3\) 的系数为 \(0\)，而 \(x\) 的系数为 \(n>0\)，不可能满足题设，所以 \(n\geq3\). 由二项式定理，\(x\) 的系数为 \(\mathrm C_n^1=n\)，\(x^3\) 的系数为 \(\mathrm C_n^3=\frac{n(n-1)(n-2)}6\). 根据题意，\(\frac{n(n-1)(n-2)}6=7n\). 因 \(n>0\)，约去 \(n\)，得 \((n-1)(n-2)=42\)，即 \(n^2-3n-40=0\). 因式分解得 \((n-8)(n+5)=0\)，所以 \(n=8\) 或 \(n=-5\). 结合 \(n\) 为正整数，得 \(n=8\).
\end{solution}

\begin{problem}
求极限：\(\displaystyle\lim_{n\to\infty}\left(\frac{1}{n^2} + \frac{2}{n^2} + \frac{3}{n^2} + \cdots + \frac{2n}{n^2}\right)\).
\end{problem}

\begin{solution}
括号内的分子依次为从 \(1\) 到 \(2n\) 的连续整数，因此 \(\frac{1}{n^2}+\frac{2}{n^2}+\cdots+\frac{2n}{n^2}=\frac1{n^2}\sum_{k=1}^{2n}k\). 利用等差数列求和公式，\(\frac1{n^2}\sum_{k=1}^{2n}k=\frac1{n^2}\cdot\frac{2n(2n+1)}2=\frac{2n^2+n}{n^2}=2+\frac1n\). 当 \(n\to\infty\) 时，\(\frac1n\to0\)，所以原极限为 \(2\).
\end{solution}

\begin{problem}
由数字 \(1\)，\(2\)，\(3\)，\(4\)，\(5\) 组成没有重复数字且数字 \(1\) 与 \(2\) 不相邻的五位数，求这种五位数的个数.
\end{problem}

\begin{solution}
五个数字各使用一次，且首位数字不可能为 \(0\)，所以不考虑相邻条件时，五位数共有 \(5!=120\) 个. 再数出 \(1\) 与 \(2\) 相邻的情形：把 \(1,2\) 看作一个整体，这个整体内部可以排成 \(12\) 或 \(21\)，有 \(2\) 种；这个整体再与 \(3\)，\(4\)，\(5\) 一起排列，共有 \(4!\) 种. 因此相邻的五位数有 \(2\cdot4!=48\) 个. 由排除法，不相邻的五位数共有 \(5!-2\cdot4!=120-48=72\)，故答案为 \(72\).
\end{solution}

\begin{problem}
求函数 \(y = \log_2(1 + 2x - 3x^2)\) 的定义域.
\end{problem}

\begin{solution}
因为对数的底数 \(2>0\) 且 \(2\ne1\)，只需令真数为正：\(1+2x-3x^2>0\). 将左边因式分解，得 \(1+2x-3x^2=-(3x+1)(x-1)\)，所以不等式等价于 \((3x+1)(x-1)<0\). 两个零点为 \(-\frac13\) 和 \(1\)，由于二次项系数为正，乘积在两个零点之间为负，故定义域为 \(\left(-\frac13,1\right)\).
\end{solution}

\begin{problem}
圆锥底面积为 \(3\pi\)，母线与底面所成的角为 \(60^\circ\)，求它的体积.
\end{problem}

\begin{solution}
设圆锥底面半径为 \(r\)，高为 \(h\)，母线长为 \(s\). 由底面积条件 \(\pi r^2=3\pi\)，得 \(r^2=3\)，即 \(r=\sqrt3\). 在圆锥的轴截面中，母线在底面上的射影是底面半径，母线与底面所成的角就是母线与该射影所成的角. 因此在轴截面的直角三角形中，\(\tan60^\circ=\frac hr\). 代入 \(r=\sqrt3\) 和 \(\tan60^\circ=\sqrt3\)，得 \(h=3\). 所以圆锥体积为 \(V=\frac13\pi r^2h=\frac13\cdot3\pi\cdot3=3\pi\).
\end{solution}

\begin{problem}
发电厂发出的电是三相交流电，它的三根导线上的电流强度分别是时间 \(t\) 的函数：\(I_A = I\sin \omega t\)，\(I_B = I\sin(\omega t + 120^\circ)\)，\(I_C = I\sin(\omega t + 240^\circ)\)，求 \(I_A + I_B + I_C\) 的值.
\end{problem}

\begin{solution}
令 \(\theta=\omega t\). 由和角公式，\(\sin(\theta+120^\circ)=-\frac12\sin\theta+\frac{\sqrt3}{2}\cos\theta\)，\(\sin(\theta+240^\circ)=-\frac12\sin\theta-\frac{\sqrt3}{2}\cos\theta\). 因此
\(\begin{aligned}
I_A+I_B+I_C
&=I\left[\sin\theta+\sin(\theta+120^\circ)+\sin(\theta+240^\circ)\right]\\
&=I\left[\sin\theta-\frac12\sin\theta+\frac{\sqrt3}{2}\cos\theta-\frac12\sin\theta-\frac{\sqrt3}{2}\cos\theta\right]=0.
\end{aligned}\)
\end{solution}

\begin{problem}
在复平面内，已知等边三角形的两个顶点表示的复数分别为 \(2\)，\(\frac{1}{2} + \frac{\sqrt{3}}{2}\i\)，求第三个顶点表示的复数.
\end{problem}

\begin{solution}
设已知两个顶点对应的复数为 \(z_1=2\)，\(z_2=\frac12+\frac{\sqrt3}{2}\i\)，并令 \(w=z_2-z_1=-\frac32+\frac{\sqrt3}{2}\i\). 等边三角形的第三个顶点到第一个顶点的向量，必须是把 \(w\) 旋转 \(60^\circ\) 或 \(-60^\circ\) 得到的向量. 因此第三个顶点对应的复数为 \(z_3=2+\left(\frac12+\frac{\sqrt3}{2}\i\right)w\) 或 \(z_3=2+\left(\frac12-\frac{\sqrt3}{2}\i\right)w\). 分别计算：
\(\left(\frac12+\frac{\sqrt3}{2}\i\right)\left(-\frac32+\frac{\sqrt3}{2}\i\right)=-\frac32-\frac{\sqrt3}{2}\i\)，所以第一种情形给出 \(z_3=2-\frac32-\frac{\sqrt3}{2}\i=\frac12-\frac{\sqrt3}{2}\i\). 又 \(\left(\frac12-\frac{\sqrt3}{2}\i\right)\left(-\frac32+\frac{\sqrt3}{2}\i\right)=\sqrt3\i\)，所以第二种情形给出 \(z_3=2+\sqrt3\i\). 故第三个顶点表示的复数为 \(\frac12-\frac{\sqrt3}{2}\i\) 或 \(2+\sqrt3\i\).
\end{solution}

\begin{problem}
如图，三棱锥 \(P-ABC\) 中，已知 \(PA \perp BC\)，\(PA = BC = L\)，\(PA\)，\(BC\) 的公垂线 \(ED = h\). 求证：三棱锥 \(P-ABC\) 的体积 \(V = \frac{1}{6}L^2h\).

\bitmapfigure[max width=0.34\linewidth]{img/1987/national_liberal/q18_fig1.png}
\end{problem}

\begin{solution}
由 \(ED\) 是 \(PA\) 与 \(BC\) 的公垂线，得 \(E\in PA\)，\(D\in BC\)，且 \(ED\perp PA\)，\(ED\perp BC\). 直线 \(PA\) 与 \(ED\) 相交并且都在平面 \(PAD\) 内；而 \(BC\) 同时垂直于这两条相交直线，所以 \(BC\perp\text{平面 }PAD\).

在平面 \(PAD\) 内，\(PA\perp ED\)，故以 \(PA\) 和 \(ED\) 为两条直角边的三角形 \(PAD\) 的面积为 \(S_{\triangle PAD}=\frac12PA\cdot ED=\frac12Lh\).

点 \(B\)，\(C\)，\(D\) 在直线 \(BC\) 上，且 \(BC\perp\text{平面 }PAD\)，所以以三角形 \(PAD\) 为底面时，三棱锥 \(B-PAD\) 和 \(C-PAD\) 的高分别为 \(BD\) 和 \(CD\). 当 \(D\) 在边 \(BC\) 上时，原三棱锥被这两个三棱锥分割，且 \(BD+CD=BC=L\). 于是
\(\begin{aligned}
V&=\frac13S_{\triangle PAD}\cdot BD+\frac13S_{\triangle PAD}\cdot CD\\
&=\frac13S_{\triangle PAD}(BD+CD)=\frac13\cdot\frac12Lh\cdot L=\frac16L^2h.
\end{aligned}\)
命题得证.
\end{solution}

\begin{problem}
设对所有实数 \(x\)，不等式 \(x^2\log_2 \frac{4(a + 1)}{a} + 2x\log_2 \frac{2a}{a + 1} + \log_2 \frac{(a + 1)^2}{4a^2} > 0\) 恒成立，求 \(a\) 的取值范围.
\end{problem}

\begin{solution}
令 \(r=\frac{2a}{a+1}\). 原不等式中各对数的真数都必须为正. 若 \(r>0\)，则 \(\frac{4(a+1)}a=\frac8r>0\)，\(\frac{(a+1)^2}{4a^2}=\frac1{r^2}>0\)，所以只需先要求 \(r>0\). 这给出 \(a>0\) 或 \(a<-1\).

令 \(u=\log_2r\)，则 \(\log_2\frac{4(a+1)}a=\log_2\frac8r=3-u\)，\(\log_2\frac{(a+1)^2}{4a^2}=\log_2\frac1{r^2}=-2u\). 原不等式化为 \(q(x)=(3-u)x^2+2ux-2u>0\) 对所有实数 \(x\) 恒成立.

若 \(3-u<0\)，则 \(x\to\infty\) 时 \(q(x)\to-\infty\)，不可能恒正；若 \(3-u=0\)，则 \(q(x)=6x-6\)，也不可能对所有 \(x\) 恒正. 因此必须 \(3-u>0\). 在此条件下，二次函数对所有实数恒正的充要条件是判别式小于零. 其判别式为 \(\Delta=(2u)^2-4(3-u)(-2u)=4u(6-u)<0\). 所以 \(u<0\) 或 \(u>6\)，与 \(u<3\) 合并得 \(u<0\). 由于对数函数 \(\log_2 r\) 的底数大于 \(1\)，\(u<0\) 等价于 \(0<r<1\).

现在解 \(0<\frac{2a}{a+1}<1\). 当 \(a>0\) 时，第二个不等式给出 \(2a<a+1\)，即 \(a<1\)，从而 \(0<a<1\). 当 \(a<-1\) 时，因分母为负，\(\frac{2a}{a+1}<1\) 等价于 \(2a>a+1\)，即 \(a>1\)，与 \(a<-1\) 矛盾. 因此 \(a\) 的取值范围为 \((0,1)\).
\end{solution}

\begin{problem}
设复数 \(z_1\) 和 \(z_2\) 满足关系式 \(z_1\overline{z_2} + \overline{A}z_1 + A\overline{z_2} = 0\)，其中 \(A\) 为不等于 \(0\) 的复数. 证明：

\begin{enumerate}
\item \(|z_1 + A||z_2 + A| = |A|^2\)；
\item \(\frac{z_1 + A}{z_2 + A} = \left|\frac{z_1 + A}{z_2 + A}\right|\).
\end{enumerate}
\end{problem}

\begin{solution}
把已知等式左边与 \((z_1+A)\overline{(z_2+A)}\) 展开式比较：\((z_1+A)\overline{(z_2+A)}=z_1\overline{z_2}+z_1\overline A+A\overline{z_2}+A\overline A\). 利用题设的前三项之和为 \(0\)，得到 \((z_1+A)\overline{(z_2+A)}=|A|^2\).

\begin{enumerate}
\item 因为 \(A\ne0\)，所以 \(|A|^2>0\). 上式右边非零，因而 \(z_2+A\ne0\)，同时 \(z_1+A\ne0\)，题中的模和与分式均有意义. 两边取模，得到 \(|z_1+A|\,|\overline{z_2+A}|=|A|^2\). 又 \(|\overline{z_2+A}|=|z_2+A|\)，所以 \(|z_1+A||z_2+A|=|A|^2\)，第 1 问得证.

\item 由 \(z_2+A\ne0\)，可将上式除以 \(|z_2+A|^2=(z_2+A)\overline{(z_2+A)}\)，得到 \(\frac{z_1+A}{z_2+A}=\frac{|A|^2}{|z_2+A|^2}>0\). 左边是正实数，正实数的模等于其本身，因此 \(\frac{z_1+A}{z_2+A}=\left|\frac{z_1+A}{z_2+A}\right|\)，第 2 问得证.
\end{enumerate}
\end{solution}

\begin{problem}
设数列 \(a_1\)，\(a_2\)，\(\cdots\)，\(a_n\)，\(\cdots\) 的前 \(n\) 项的和 \(S_n\) 与 \(a_n\) 满足 \(S_n = ka_n + 1\)（其中 \(k\) 是与 \(n\) 无关的常数，且 \(k \neq 1\)).

\begin{enumerate}
\item 试写出用 \(n\)，\(k\) 表示的 \(a_n\) 的表达式；
\item 若 \(\displaystyle\lim_{n\to\infty}S_n = 1\)，求 \(k\) 的取值范围.
\end{enumerate}
\end{problem}

\begin{solution}
\begin{enumerate}
\item 取 \(n=1\)，由 \(S_1=a_1=ka_1+1\) 得 \(a_1=\frac1{1-k}\). 对 \(n\geq1\)，利用 \(S_{n+1}-S_n=a_{n+1}\) 以及题设关系，得 \(a_{n+1}=k(a_{n+1}-a_n)\)，即 \((k-1)a_{n+1}=ka_n\).

当 \(k\ne0\) 时，\(a_1\ne0\)，且递推式给出 \(\frac{a_{n+1}}{a_n}=\frac{k}{k-1}\). 因此 \(a_n=a_1\left(\frac{k}{k-1}\right)^{n-1}=\frac1{1-k}\left(\frac{k}{k-1}\right)^{n-1}\). 当 \(k=0\) 时，题设直接给出 \(S_n=1\)，所以 \(a_1=S_1=1\)，并且 \(a_n=S_n-S_{n-1}=0\)（\(n\geq2\)）.

\item 由 \(S_n-1=ka_n\)，所求条件等价于 \(ka_n\to0\). 当 \(k=0\) 时，\(S_n=1\) 对所有 \(n\) 成立，条件满足. 当 \(k\ne0\) 时，由上面的通项式，\(ka_n\) 是非零常数乘以 \(\left(\frac{k}{k-1}\right)^{n-1}\). 该幂数列趋于 \(0\) 的充要条件为 \(\left|\frac{k}{k-1}\right|<1\). 因 \(k\) 为实数且 \(k\ne1\)，两边平方得 \(k^2<(k-1)^2=k^2-2k+1\)，即 \(2k<1\)，所以 \(k<\frac12\). 其中已包含 \(k=0\)，且不会包含被题设排除的 \(k=1\)，故最终范围为 \(k<\frac12\).
\end{enumerate}
\end{solution}

\begin{problem}
正方形 \(ABCD\) 在直角坐标平面内，已知其一条边 \(AB\) 在直线 \(y = x + 4\) 上，\(C\)，\(D\) 在抛物线 \(x = y^2\) 上，求正方形 \(ABCD\) 的面积.
\end{problem}

\begin{solution}
正方形的对边平行，所以 \(CD\parallel AB\). 设直线 \(CD\) 的方程为 \(y=x+b\)，即 \(x=y-b\). 设 \(C\)，\(D\) 的纵坐标分别为 \(y_1\)，\(y_2\). 因为 \(C\)，\(D\) 在抛物线 \(x=y^2\) 和直线 \(x=y-b\) 上，所以 \(y_1\)，\(y_2\) 是方程 \(y^2=y-b\)，即 \(y^2-y+b=0\) 的两个根. 由韦达定理，\(y_1+y_2=1\)，\(y_1y_2=b\)，从而 \((y_1-y_2)^2=(y_1+y_2)^2-4y_1y_2=1-4b\).

对应点的横坐标为 \(x_i=y_i-b\)，所以 \(x_1-x_2=y_1-y_2\). 因此 \(CD^2=(x_1-x_2)^2+(y_1-y_2)^2=2(1-4b)=2-8b\).

两条平行直线 \(AB:y=x+4\) 与 \(CD:y=x+b\) 的距离为 \(\frac{|4-b|}{\sqrt2}\). 这正是正方形的边长，因为它等于两条对边之间的距离；又 \(AB=CD\)，故 \(\frac{(4-b)^2}{2}=CD^2=2-8b\). 化简得 \((4-b)^2=4-16b\)，即 \(b^2+8b+12=0\)，所以 \(b=-2\) 或 \(b=-6\). 两者都使 \(1-4b>0\)，因而直线 \(CD\) 都与抛物线有两个不同交点，确实能形成相应的正方形.

当 \(b=-2\) 时，正方形面积为 \(CD^2=2-8(-2)=18\)；当 \(b=-6\) 时，正方形面积为 \(CD^2=2-8(-6)=50\). 故正方形的面积为 \(18\) 或 \(50\).
\end{solution}
